% Nota corta — frontera cúbica de Erdős–Gyárfás mediante SMS calibrado.
% Los campos [[...]] se rellenan antes de publicar.
\documentclass[11pt]{article}
\usepackage[utf8]{inputenc}
\usepackage{amsmath,amsthm,amssymb}
\usepackage{hyperref}
\usepackage{booktabs}
\usepackage{microtype}

\newtheorem{theorem}{Theorem}
\newtheorem{conjecture}{Conjecture}
\newtheorem{proposition}{Proposition}

\title{Cubic graphs on at most 40 vertices satisfy the
  Erd\H{o}s--Gy\'arf\'as conjecture:\\
  a calibrated SAT-modulo-symmetries search}
\author{Christopher Andre Manzano Vimos\\
  \small WhiteBox ML, Madrid\\
  \small \texttt{christopher.manzano@manapple.dev} \quad
  ORCID: \href{https://orcid.org/0009-0000-4270-4610}{0009-0000-4270-4610}}
\date{August 2026}

\begin{document}
\maketitle

\begin{abstract}
The Erd\H{o}s--Gy\'arf\'as conjecture states that every graph with minimum
degree at least~3 contains a simple cycle whose length is a power of two.
Markstr\"om (2004) verified the conjecture for cubic graphs on at most 28
vertices, and that exhaustive bound stood until 2026. We report an exhaustive
verification for connected cubic graphs on 30, 32, 34, 36, 38 and 40 vertices:
in each
order there is no cubic graph avoiding cycles of length 4, 8 and 16.
Consequently a smallest cubic counterexample, if one exists, has at least 42
vertices. The search uses SAT modulo symmetries with a complete
forbidden-subgraph propagator; unlike a decision procedure taken on trust, the
instrument is first \emph{calibrated} against exhaustive enumeration: it is
required to reproduce, exactly and as sets, thirteen reference censuses of cubic
graphs, nine of which have nonzero answers, including a 348.7 core-hour
enumeration of order~30 carried out independently by the author. Positive
controls are additionally run at the frontier orders themselves. All code,
logs and censuses are available for independent verification.
\end{abstract}

\section{Introduction}

Erd\H{o}s and Gy\'arf\'as asked whether every finite graph with minimum degree
at least~3 contains a cycle of length $2^k$ for some $k\ge 2$; the problem
appears in Erd\H{o}s's problem papers from 1993 onwards~\cite{erdos93}, and
both believed the answer to be negative. Erd\H{o}s offered \$100 for a proof
and \$50 for a counterexample~\cite{markstrom2004}; the problem is listed with
a \$1000 prize as \#64 of the Erd\H{o}s problems
collection~\cite{erdosproblems}. Liu and Montgomery~\cite{liumontgomery}
settled the case of sufficiently large minimum degree, and
Carr~\cite{carr2026} showed that any minimal counterexample is predominantly
cubic, so the cubic case is the structurally critical one.

\paragraph{Computational state of the art.}
Markstr\"om~\cite{markstrom2004} generated all cubic graphs on fewer than 29
vertices and found no counterexample; his Table~3 records the numbers of cubic
graphs with neither a 4-cycle nor an 8-cycle for $n=24,26,28$, namely $4$, $23$
and $251$. In 2026 Balaji~\cite{balaji2026} reported, via SAT modulo
symmetries, that no graph of minimum degree at least~3 on at most 31 vertices
avoids all power-of-two cycles, which yields the bound $n\ge 32$ for the
general class and, a fortiori, for the cubic one; that work is a preprint under
review, and its documentation states that no machine-checkable certificate is
currently available for the claim, since the clauses contributed by the
subgraph propagator are not RUP/RAT-derivable from the encoding, and that
third-party reproduction is outstanding. Tranquilli~\cite{tranquilli2026}
treats the cubic \emph{bipartite} subclass and reports the bound $n\ge 60$
there.

\paragraph{Result.}
\begin{theorem}\label{thm:main}
Every connected cubic graph on $n \le 40$ vertices contains a cycle of length
4, 8 or 16. In particular, a smallest cubic counterexample to the
Erd\H{o}s--Gy\'arf\'as conjecture has at least 42 vertices.
\end{theorem}

Cubic graphs have even order, and a counterexample may be assumed connected,
since every component of a cubic graph is itself cubic. Orders $n\le 28$ are
due to Markstr\"om; orders 30, 32, 34, 36, 38 and 40 are verified here.

\paragraph{Why forbidding $\{4,8,16\}$ suffices.}
For $n\le 31$ the admissible power-of-two cycle lengths are exactly
$\{4,8,16\}$, so avoiding them is equivalent to being a counterexample. From
$n=32$ on, $C_{32}$ also fits, and adding a 32-cycle to the forbidden set
inside the search is expensive: at $n=32$ it is a Hamiltonian cycle, and as a
propagator it would be evaluated on every partial assignment. We therefore
\emph{decouple}: the search forbids only $\{4,8,16\}$, which is a
\emph{weaker} constraint, and the larger powers of two are tested afterwards
on whatever survives. Since the search returns no graph at all in every order
considered, the posterior test is vacuous and the conclusion is stronger than
required: not merely ``no counterexample'', but ``every cubic graph of that
order already contains one of $C_4$, $C_8$, $C_{16}$''.

\section{Method}

We use SAT modulo symmetries (SMS, Kirchweger and Szeider~\cite{sms}), which
performs complete, isomorph-free graph generation inside a CDCL solver, built
together with the Glasgow subgraph solver~\cite{glasgow} as a complete
forbidden-subgraph propagator. One call per order decides: \emph{is there a
connected cubic graph on $n$ vertices containing no $C_4$, $C_8$ or $C_{16}$
as a subgraph?} The graph is encoded by the $\binom n2$ edge variables; cubicity
is imposed as a pair of cardinality constraints (sequential counter) fixing
every degree to exactly~3, and connectivity is imposed by SMS.

Two configuration choices deserve mention. First, the minimality (canonicity)
check of SMS is applied without cutoff; the default cutoff of $2\cdot10^5$
truncates the check and thereby weakens complete symmetry breaking. Second, a
run that exhausts its time budget is recorded as such and never as
unsatisfiability, and a run whose solution count cannot be parsed is recorded
as an error rather than as zero; exit codes are checked in every case.

\subsection{Calibration of the instrument}

An unsatisfiability answer is an assertion of the solver, not an object one can
inspect; enumeration, by contrast, produces countable objects that can be
compared with published censuses. The dangerous failure mode of a
SAT-plus-propagator pipeline is over-constraint: such a pipeline returns zero
everywhere and therefore passes every test whose correct answer is zero. Only a
test with a known \emph{nonzero} answer detects it.

We therefore require the pipeline to reproduce thirteen reference censuses of
connected cubic graphs, nine of them nonzero (Table~\ref{tab:calib}). Three
sources of ground truth are used: Markstr\"om's published Table~3; an
exhaustive enumeration at orders 24--30 performed by the author with
\texttt{geng} (nauty~\cite{nauty}) and an embedded C pruning hook, 348.7
core-hours at order~30; and, to exercise the $C_{16}$ propagator against
nonzero answers, fresh counts of cubic graphs containing no 16-cycle at orders
14--22. At orders 14--20 these counts were computed twice with independent
cycle detectors (a bitmask depth-first search and a networkx-based one). Order
22, the largest, was counted once, with the bitmask detector, and guarded
differently: it was generated in six \texttt{res/mod} shards whose graph totals
sum to $7\,319\,447$, the known number of connected cubic graphs on 22
vertices, so the partition neither lost nor duplicated a graph; of those,
$52\,781$ contain no 16-cycle, and each of those $52\,781$ was re-tested
individually with the networkx detector. The agreement of SMS at this order is
not counted as verification of the anchor, since SMS is the instrument under
calibration.

\begin{table}[h]\centering
\begin{tabular}{llll}
\toprule
Order & Forbidden & Expected & SMS \\
\midrule
$n=10$ & $C_{16}$ & 19 & 19 \\
$n=14$ & $C_{16}$ & 509 (no 16-cycle fits) & 509 \\
$n=16$ & $C_{16}$ & 219 & 219 \\
$n=18$ & $C_{16}$ & 1471 & 1471 \\
$n=20$ & $C_{16}$ & 12709 & 12709 \\
$n=22$ & $C_{16}$ & 52781 & 52781 \\
$n=24$ & $C_4,C_8$ & 4 (Markstr\"om) & 4, same graphs \\
$n=26$ & $C_4,C_8$ & 23 (Markstr\"om) & 23, same graphs \\
$n=28$ & $C_4,C_8$ & 251 (Markstr\"om) & 251 \\
$n=24,26,28,30$ & $C_4,C_8,C_{16}$ & 0 & 0 \\
\bottomrule
\end{tabular}
\caption{Calibration against exhaustive enumeration. At $n=24$ and $n=26$ the
comparison is at the level of sets, not merely cardinalities: the graphs
emitted by SMS and by \texttt{geng} coincide after canonical labelling with
\texttt{labelg}.}\label{tab:calib}
\end{table}

Two robustness axes were varied. Re-deciding with the totalizer cardinality
encoding, a structurally different CNF for the same degree constraint,
reproduces $251$ at $n=28$, $1471$ at $n=18$ and zero at $n=30$, $n=32$,
$n=36$ and $n=40$; the last two are frontier orders, re-decided in 4143\,s and
17\,207\,s against 3366\,s and 18\,922\,s for the sequential counter, so the
largest order reported here has been decided twice with structurally different
formulas. The comparison of graph sets at $n=24$ and $n=26$ is exact.

\subsection{Positive controls at the frontier orders}

Calibration establishes that the encoding and the three propagators are
correct, and encoding faults are order-independent: the same code runs at every
order. It does not, however, exclude order-dependent breakage that would render
the formula vacuously unsatisfiable above some size. We therefore also run
existence queries at the frontier orders themselves, forbidding only
$\{C_4,C_8\}$, where solutions are known to exist: SMS returns a satisfying
graph at $n=32$, $34$, $36$, $38$, $40$, $42$ and $44$, the last two beyond the
orders decided here. The zero counts reported below are
thus not an artifact of the pipeline ceasing to find anything at large orders.

\subsection{Results}

\begin{table}[h]\centering
\begin{tabular}{lll}
\toprule
$n$ & connected cubic graphs with no $C_4,C_8,C_{16}$ & single-core time \\
\midrule
30 & 0 & 187\,s \\
32 & 0 & 505\,s \\
34 & 0 & 1188\,s \\
36 & 0 & 3366\,s \\
38 & 0 & 8157\,s \\
40 & 0 & 18922\,s \\
\bottomrule
\end{tabular}
\caption{The frontier sweep. Order~30 also serves as calibration: the same
answer was obtained by exhaustive enumeration in 348.7 core-hours.}
\label{tab:frontier}
\end{table}

The observed growth is a factor of roughly $2.3$ to $2.8$ per two vertices
(order~40 came in at $2.32$ times order~38), so the method remains within reach
of a single desktop machine for several further orders. Order~40 was decided in
a single-threaded run of 5\,h\,15\,min; the native cube-and-conquer
parallelisation of SMS was measured on this instance family and discarded, as
it yields only a factor of~$1.7$ on twelve workers, and CaDiCaL's formula
simplification, which is faster still, does not preserve the set of models:
under it the $n=28$ census returns 241 instead of 251.

\section{What is not claimed}

The result above order~30 rests on a single method. No machine-checkable proof
certificate accompanies it: SMS can emit an LRAT proof, but the clauses
contributed by the subgraph propagator are not RUP/RAT-derivable from the
encoding alone, so a generic checker cannot validate it; a certified-SMS
treatment in the sense of~\cite{sms} is the natural strengthening and is not
attempted here. What we do claim is a calibrated instrument: the pipeline
reproduces every reference census available in this class, including eight with
nonzero answers and two at the level of graph sets, and demonstrates
satisfiable behaviour live at the frontier orders. Order-dependent failure
modes are guarded procedurally, not proved absent. Independent reproduction
remains desirable, and the artifacts below are published to that end.

\section{Data availability}

Repository:
\url{https://github.com/christopher-manzanovimos_geci/erdos-gyarfas-cubic-frontier};
archived deposit: Zenodo, DOI
\href{https://doi.org/10.5281/zenodo.21976755}{10.5281/zenodo.21976755}.
The repository contains the
SMS driver, the calibration battery, the ground-truth generators and cycle
detectors, the run logs with exit codes and timings, the graph6 censuses for
$n=14$--$22$ (no $C_{16}$) and $n=24$, $26$ (no $C_4/C_8$), the complete
July 2026 enumeration apparatus used as independent ground truth, and the build
record. Versions used: SMS at commit \texttt{464f12f}, CaDiCaL at
\texttt{b023aaf}, Glasgow subgraph solver at \texttt{1217f5b}; note that the
SMS build script clones the Glasgow solver at its current head rather than at
a pinned commit, so this component should be recorded rather than assumed by
anyone reproducing the computation.

\section*{Acknowledgements}
The author thanks WhiteBox ML for supporting the publication of this work. The
computational pipeline was developed with the
assistance of Claude (Anthropic); the author is responsible for the design of
the experiments, the verification protocol and all claims made here.

\begin{thebibliography}{9}
\bibitem{erdos93} P.~Erd\H{o}s, Some of my favorite solved and unsolved
  problems in graph theory, \emph{Quaestiones Math.} 16 (1993) 333--350.
\bibitem{erdosproblems} T.~F. Bloom, Erd\H{o}s Problems, \#64,
  \url{https://www.erdosproblems.com/64} (accessed August 2026).
\bibitem{markstrom2004} K.~Markstr\"om, Extremal graphs for some problems on
  cycles in graphs, \emph{Congr. Numer.} 171 (2004) 179--192.
\bibitem{liumontgomery} H.~Liu and R.~Montgomery, A solution to Erd\H{o}s and
  Hajnal's odd cycle problem, \emph{J. Amer. Math. Soc.} 36 (2023) 1191--1234.
\bibitem{carr2026} A.~Carr, Every minimal counterexample to the
  Erd\H{o}s--Gy\'arf\'as conjecture is predominantly cubic, arXiv:2605.22844
  (2026).
\bibitem{balaji2026} A.~Balaji, Verifying the Erd\H{o}s--Gy\'arf\'as conjecture
  up to 31 vertices with SAT modulo symmetries, Zenodo (2026),
  \url{https://doi.org/10.5281/zenodo.21190438}; accompanying repository
  \url{https://github.com/ArjunBalaji79/erdos-gyarfas-min-degree-3}
  (commit \texttt{53502b0f}, 3 July 2026; accessed August 2026).
\bibitem{tranquilli2026} J.~Tranquilli, A 60-vertex lower bound for cubic
  bipartite counterexamples to the Erd\H{o}s--Gy\'arf\'as conjecture,
  arXiv:2608.02675 (2026).
\bibitem{sms} M.~Kirchweger and S.~Szeider, SAT modulo symmetries for graph
  generation and enumeration, \emph{ACM Trans. Comput. Logic} 25 (2024),
  Article~18.
\bibitem{glasgow} C.~McCreesh, P.~Prosser and J.~Trimble, The Glasgow subgraph
  solver: using constraint programming to tackle hard subgraph isomorphism
  problem variants, in \emph{ICGT 2020}, LNCS 12150, 316--324.
\bibitem{nauty} B.~D. McKay and A.~Piperno, Practical graph isomorphism, II,
  \emph{J. Symbolic Comput.} 60 (2014) 94--112.
\end{thebibliography}

\end{document}
